% Volume VII: Deep Learning Mathematics
% Continuity note for Chapter 41.
% Previous dependencies: Chapters 13-18 introduced probability, densities, Bayes' theorem, KL divergence, cross-entropy, and mutual information; Chapters 22 and 34 introduced variational principles and latent variables; Chapters 27 and 40 provide stochastic-process and architecture background.
% New durable concepts: latent-variable generation, inference as posterior reversal, ELBOs, reparameterization, invertible flows and change of variables, adversarial games, energy-based models, partition functions, diffusion noising, score functions, and reverse SDE intuition.
% Forward hooks: Chapter 42 studies representation geometry and robustness of learned generative features; Chapters 49-50 reuse energy, free energy, and inference ideas for neural circuits and predictive coding.
\section{Chapter 41: Generative Deep Learning}
\label{ch:generative-deep-learning}

Discriminative models learn maps such as \(x\mapsto y\). Generative models learn how data could have been produced. They specify, approximate, or simulate a probability distribution over observations. The generative direction usually runs from hidden causes to data: a latent variable, noise source, or intermediate state produces an observation. Inference reverses that direction: after observing \(x\), we ask which hidden causes or denoising steps could explain it.

\paragraph{Chapter problem.}
How do deep generative models represent probability distributions over complex data, and where do Bayes' theorem, KL divergence, change of variables, adversarial games, energy normalization, and stochastic differential equations enter?

\paragraph{Section map.}
Section 41.1 derives variational autoencoders as latent-variable models trained by an evidence lower bound. Section 41.2 studies normalizing flows as invertible maps trained by exact likelihood through change of variables. Section 41.3 studies GANs as adversarial distribution matching and introduces energy-based models as a bridge to statistical physics. Section 41.4 studies diffusion and score-based models as learned reversals of gradual noising processes.

\subsection{41.1 Variational autoencoders}

\paragraph{Guiding question.}
How can a neural network learn a latent-cause model \(z\to x\) when the posterior \(p(z\mid x)\) is hard to compute exactly?

\paragraph{Beginner-facing intuition.}
A variational autoencoder (VAE) assumes each observation \(x\) was generated from an unobserved latent variable \(z\). The decoder \(p_\theta(x\mid z)\) maps latent causes to data distributions. The prior \(p(z)\) says which latent causes are plausible before observing data. After observing \(x\), the posterior \(p_\theta(z\mid x)\) tells us which latent causes explain it.

The problem is that the exact posterior usually requires an intractable integral. A VAE introduces an encoder \(q_\phi(z\mid x)\) to approximate that posterior. Training balances two pressures: reconstruct the observation well from sampled \(z\), and keep the approximate posterior close to the prior so that new samples from the prior decode into realistic data.

\paragraph{Formal definitions and notation.}
Let \(X\) be an observation in a data space \(\mathcal X\), and let \(Z\in\mathbb{R}^{d_z}\) be a latent variable. A latent-variable generative model is
\[
p_\theta(x,z)=p(z)p_\theta(x\mid z).
\]
The marginal likelihood of \(x\) is
\[
p_\theta(x)=\int p(z)p_\theta(x\mid z)\,dz,
\]
and Bayes' theorem gives the posterior
\[
p_\theta(z\mid x)=\frac{p(z)p_\theta(x\mid z)}{p_\theta(x)}.
\]
A VAE uses an approximate posterior, or encoder,
\[
q_\phi(z\mid x),
\]
and trains \(\theta,\phi\) by maximizing the evidence lower bound (ELBO)
\[
\boxed{
\mathcal L_{\mathrm{ELBO}}(x)
=
\mathbb{E}_{q_\phi(z\mid x)}[\log p_\theta(x\mid z)]
-D_{\mathrm{KL}}(q_\phi(z\mid x)\|p(z)).
}
\]

\paragraph{Core derivation: ELBO plus posterior KL equals log evidence.}
Start with any density \(q_\phi(z\mid x)\) whose support is compatible with the posterior. Then
\[
\begin{aligned}
\log p_\theta(x)
&=\log p_\theta(x)\int q_\phi(z\mid x)\,dz\\
&=\mathbb{E}_{q_\phi}\left[\log\frac{p_\theta(x,z)}{q_\phi(z\mid x)}\right]
+\mathbb{E}_{q_\phi}\left[\log\frac{q_\phi(z\mid x)}{p_\theta(z\mid x)}\right]\\
&=\mathcal L_{\mathrm{ELBO}}(x)
+D_{\mathrm{KL}}(q_\phi(z\mid x)\|p_\theta(z\mid x)).
\end{aligned}
\]
Since KL divergence is nonnegative,
\[
\mathcal L_{\mathrm{ELBO}}(x)\leq \log p_\theta(x).
\]
Maximizing the ELBO therefore increases a lower bound on log likelihood while also pushing the encoder toward the exact posterior.

For a Gaussian encoder
\[
q_\phi(z\mid x)=\mathcal N(\mu_\phi(x),\operatorname{diag}(\sigma_\phi^2(x))),
\]
the reparameterization trick writes
\[
z=\mu_\phi(x)+\sigma_\phi(x)\odot\epsilon,\qquad
\epsilon\sim\mathcal N(0,I),
\]
so gradients can flow through \(\mu_\phi\) and \(\sigma_\phi\).

\paragraph{Concrete example.}
Let \(p(z)=\mathcal N(0,1)\) and \(q(z\mid x)=\mathcal N(\mu,\sigma^2)\). The KL term has the closed form
\[
D_{\mathrm{KL}}(\mathcal N(\mu,\sigma^2)\|\mathcal N(0,1))
=\frac12(\mu^2+\sigma^2-\log\sigma^2-1).
\]
If \(\mu=1\) and \(\sigma^2=1\), then the KL is \(1/2\). If \(\mu=0\) and \(\sigma^2=1\), the KL is \(0\). Thus the KL term penalizes moving the approximate posterior away from the prior unless doing so improves reconstruction enough.

\paragraph{Computational neuroscience example.}
A VAE can model neural population activity by positing low-dimensional latent variables \(z\) representing task state, movement, arousal, or other hidden factors, with \(p_\theta(x\mid z)\) generating observed spike counts or calcium activity. The encoder estimates latent state from observed activity. This is a useful statistical model, but the latent variables should not be declared biological causes without experimental validation.

\paragraph{AI/ML example.}
In image modeling, a VAE decoder maps a latent vector to a distribution over pixels. Sampling proceeds in the generative direction:
\[
z\sim p(z),\qquad x\sim p_\theta(x\mid z).
\]
Encoding an existing image runs the approximate inference direction:
\[
x\mapsto q_\phi(z\mid x).
\]
This separation between generation and inference is the central probabilistic idea.

\paragraph{Common misconceptions and failure modes.}
The encoder is not the generative model; it is an inference approximation. The ELBO is not exactly the log likelihood unless \(q_\phi(z\mid x)=p_\theta(z\mid x)\). A strong decoder can ignore \(z\), producing posterior collapse. A smooth latent space is encouraged by the KL term, not guaranteed by the word ``autoencoder.'' A low-dimensional latent variable is interpretable only if the model and data support that interpretation.

\paragraph{Exercises.}
\begin{enumerate}[label=\textbf{41.1.\arabic*.}, leftmargin=*]
\item \textbf{Mechanical.} Write the joint density \(p_\theta(x,z)\), marginal likelihood \(p_\theta(x)\), and posterior \(p_\theta(z\mid x)\) for a latent-variable model with prior \(p(z)\) and decoder \(p_\theta(x\mid z)\).
\item \textbf{Proof/derivation.} Derive \(\log p_\theta(x)=\mathcal L_{\mathrm{ELBO}}(x)+D_{\mathrm{KL}}(q_\phi(z\mid x)\|p_\theta(z\mid x))\).
\item \textbf{Numerical/computational.} Compute the Gaussian KL term for \(\mu=0.5\) and \(\sigma^2=2\) using the formula above.
\item \textbf{Interpretation.} In a neural VAE, what is the difference between a latent variable that predicts population activity and a latent variable that has been experimentally established as a biological cause?
\item \textbf{Misconception/debugging.} A student says, ``A VAE reconstructs data, so it is just PCA with nonlinear layers.'' Explain the missing probabilistic pieces: prior, likelihood, posterior approximation, and KL term.
\end{enumerate}

\subsection{41.2 Normalizing flows}

\paragraph{Guiding question.}
How can an invertible neural network turn a simple base distribution into a complex data distribution while keeping exact likelihoods tractable?

\paragraph{Beginner-facing intuition.}
A normalizing flow starts with a simple random variable \(Z\), often standard Gaussian, and transforms it through an invertible map into data \(X\). If the map stretches space, density decreases; if it compresses space, density increases. Change of variables tells us exactly how density changes.

Flows trade architectural freedom for exact density evaluation. Every layer must be invertible, and its Jacobian determinant must be computable. This constraint is why coupling layers, autoregressive flows, and invertible \(1\times1\) convolutions are designed carefully.

\paragraph{Formal definitions and notation.}
Let \(f_\theta:\mathcal X\to\mathcal Z\) be an invertible differentiable map with inverse \(g_\theta=f_\theta^{-1}\). We often write
\[
z=f_\theta(x),\qquad x=g_\theta(z),
\]
where \(p_Z(z)\) is a simple base density. The change-of-variables formula gives
\[
\boxed{
\log p_X(x)
=\log p_Z(f_\theta(x))
+\log\left|\det J_{f_\theta}(x)\right|,
}
\]
where \(J_{f_\theta}(x)=\partial f_\theta(x)/\partial x\). Equivalently, if \(x=g_\theta(z)\),
\[
\log p_X(g_\theta(z))
=\log p_Z(z)-\log|\det J_{g_\theta}(z)|.
\]
Training by maximum likelihood minimizes the negative log likelihood of observed data under \(p_X\).

\paragraph{Core derivation: density changes by inverse volume.}
For a one-dimensional increasing transformation \(z=f(x)\), probability mass is preserved:
\[
p_X(x)\,dx\approx p_Z(z)\,dz.
\]
Since \(dz=f'(x)\,dx\),
\[
p_X(x)=p_Z(f(x))|f'(x)|.
\]
In \(d\) dimensions, the derivative becomes the Jacobian matrix, and local volume scales by \(|\det J_f(x)|\). This yields the multidimensional formula above. The determinant is the mathematical price of exact likelihood.

\paragraph{Concrete example.}
Let \(Z\sim\mathcal N(0,1)\) and define \(z=f(x)=(x-1)/2\). Then \(x=2z+1\), and
\[
\left|\frac{dz}{dx}\right|=\frac12.
\]
Thus
\[
p_X(x)=p_Z\left(\frac{x-1}{2}\right)\frac12.
\]
At \(x=1\), \(z=0\), so
\[
p_X(1)=\frac{1}{\sqrt{2\pi}}\cdot\frac12.
\]
The factor \(1/2\) appears because the map \(x=2z+1\) stretches the base density by a factor of \(2\).

\paragraph{Computational neuroscience example.}
Flows can model complicated distributions of neural population activity by transforming a simple latent density into a flexible response density. Because likelihood is exact, one can compare held-out log likelihoods across models. The invertibility constraint, however, means a basic flow has the same dimension for \(x\) and \(z\), which may be awkward when the scientific hypothesis involves lower-dimensional latent causes.

\paragraph{AI/ML example.}
In image generation, flows can sample by drawing \(z\sim\mathcal N(0,I)\) and computing \(x=g_\theta(z)\). They can also evaluate exact log likelihood for an image by computing \(z=f_\theta(x)\) and the log determinant. This distinguishes them from GANs, which sample easily but do not provide a normalized likelihood by default.

\paragraph{Common misconceptions and failure modes.}
Invertible does not mean interpretable. A flow's latent coordinates may be convenient for density modeling without aligning with semantic factors. Exact likelihood does not guarantee perceptual sample quality. The determinant can be expensive unless the architecture is designed for tractability. A flow cannot directly compress dimension without adding noise, augmentation, or a different modeling strategy.

\paragraph{Exercises.}
\begin{enumerate}[label=\textbf{41.2.\arabic*.}, leftmargin=*]
\item \textbf{Mechanical.} For \(z=3x\) and \(Z\sim\mathcal N(0,1)\), write \(p_X(x)\) using change of variables.
\item \textbf{Proof/derivation.} Starting from preservation of probability mass in one dimension, derive \(p_X(x)=p_Z(f(x))|f'(x)|\).
\item \textbf{Numerical/computational.} If \(f(x)=x/4\), compute \(\log|\det J_f|\) in one dimension.
\item \textbf{Interpretation.} Why does exact likelihood in a flow require both invertibility and a tractable Jacobian determinant?
\item \textbf{Misconception/debugging.} A student proposes a flow with a ReLU layer \(z=\max(0,Wx)\). Explain why this is generally not a valid normalizing-flow layer.
\end{enumerate}

\subsection{41.3 GANs}

\paragraph{Guiding question.}
How can a model learn to generate realistic samples when it does not explicitly compute a likelihood?

\paragraph{Beginner-facing intuition.}
A generative adversarial network (GAN) contains two models. The generator transforms noise into fake samples. The discriminator tries to distinguish real data from generated data. Training is a game: the discriminator improves its test, and the generator improves its samples to fool the discriminator.

GANs are attractive because sampling is direct and the generator can be very flexible. Their difficulty is that the objective is a moving two-player game, not a single stable likelihood. Mode collapse, oscillation, and sensitivity to architecture or learning rates are common.

\paragraph{Formal definitions and notation.}
Let \(Z\sim p_Z\) be noise, and let
\[
G_\theta:\mathcal Z\to\mathcal X
\]
be a generator. It induces a model distribution \(p_g\) on samples \(X=G_\theta(Z)\). The discriminator
\[
D_\psi:\mathcal X\to(0,1)
\]
estimates the probability that an input is real. The original GAN objective is
\[
\boxed{
\min_\theta\max_\psi
\left[
\mathbb{E}_{x\sim p_{\mathrm{data}}}\log D_\psi(x)
+
\mathbb{E}_{z\sim p_Z}\log(1-D_\psi(G_\theta(z)))
\right].
}
\]
The generator maps latent noise to observations; there is no encoder or exact posterior unless an additional model is introduced.

\paragraph{Core derivation: the optimal discriminator for fixed generator.}
For fixed \(p_g\), the discriminator objective can be written pointwise as
\[
\int \left[p_{\mathrm{data}}(x)\log D(x)+p_g(x)\log(1-D(x))\right]\,dx.
\]
For each \(x\), maximize
\[
a\log D+b\log(1-D),
\]
where \(a=p_{\mathrm{data}}(x)\) and \(b=p_g(x)\). Differentiating gives
\[
\frac{a}{D}-\frac{b}{1-D}=0,
\]
so
\[
\boxed{
D^\star(x)=\frac{p_{\mathrm{data}}(x)}
{p_{\mathrm{data}}(x)+p_g(x)}.
}
\]
If \(p_g=p_{\mathrm{data}}\), then \(D^\star(x)=1/2\): the discriminator cannot do better than chance. Substituting \(D^\star\) connects the idealized GAN objective to Jensen-Shannon divergence, but practical training uses finite networks and stochastic updates.

\paragraph{Energy-based bridge.}
An energy-based model defines
\[
p_\theta(x)=\frac{\exp(-E_\theta(x))}{Z_\theta},
\qquad
Z_\theta=\int \exp(-E_\theta(x))\,dx.
\]
Low energy means high probability after normalization. The local learning mechanism appears by differentiating
\[
\log p_\theta(x)=-E_\theta(x)-\log Z_\theta.
\]
Assuming the derivative can pass through the integral,
\[
\nabla_\theta \log Z_\theta
=\frac{1}{Z_\theta}\nabla_\theta\int \exp(-E_\theta(u))\,du
=-\mathbb{E}_{X\sim p_\theta}\!\left[\nabla_\theta E_\theta(X)\right],
\]
so
\[
\boxed{
\nabla_\theta \log p_\theta(x)
=-\nabla_\theta E_\theta(x)
+\mathbb{E}_{X\sim p_\theta}\!\left[\nabla_\theta E_\theta(X)\right].
}
\]
For maximum likelihood, the data term is the positive phase: it changes parameters to lower energy on observed samples. The model expectation is the negative phase: it raises energy on samples the current model itself tends to produce. The partition function \(Z_\theta\) is difficult because the model expectation ranges over the whole data space, so practical EBM learning usually needs sampling or approximation, such as Markov chain Monte Carlo, contrastive divergence, score matching, or noise-contrastive estimation. This is the bridge to statistical physics: \(E_\theta\) is an energy and \(Z_\theta\) is a partition function. Unlike a GAN, an energy model defines an unnormalized density; unlike a flow, its normalization is usually intractable.

\paragraph{Concrete example.}
Suppose at one point \(x\), the data density is \(p_{\mathrm{data}}(x)=0.6\) and the generator density is \(p_g(x)=0.2\). The optimal discriminator value is
\[
D^\star(x)=\frac{0.6}{0.6+0.2}=0.75.
\]
At points where the generator puts too little mass compared with data, the optimal discriminator is confident that samples are real. The generator receives pressure, through the discriminator's gradient, to increase mass in such regions.

\paragraph{Computational neuroscience example.}
Adversarial objectives can be used to generate synthetic neural activity that matches selected statistics of recorded activity. Energy-based models also appear in simplified attractor and Hopfield-style models, where lower energy states are more stable or more probable. These links are mathematical analogies unless the energy function is tied to a concrete biophysical mechanism.

\paragraph{AI/ML example.}
GANs have been used for image synthesis, super-resolution, and data augmentation. A generator can produce sharp samples because it is trained through a discriminator sensitive to sample realism. However, the lack of explicit likelihood makes coverage hard to measure, and mode collapse can produce realistic but insufficiently diverse samples.

\paragraph{Common misconceptions and failure modes.}
The discriminator is not a permanent judge of truth; it changes during training and can overfit. Fooling a discriminator does not guarantee matching the full data distribution. Mode collapse means the generator covers only part of the data distribution. Energy values are not probabilities until normalized by the partition function. A low energy for data is insufficient if the model assigns low energy to too many non-data states.

\paragraph{Exercises.}
\begin{enumerate}[label=\textbf{41.3.\arabic*.}, leftmargin=*]
\item \textbf{Mechanical.} Write the GAN generator sampling procedure from \(z\sim p_Z\) to \(x=G_\theta(z)\).
\item \textbf{Proof/derivation.} Derive \(D^\star(x)=p_{\mathrm{data}}(x)/(p_{\mathrm{data}}(x)+p_g(x))\) for a fixed generator.
\item \textbf{Numerical/computational.} Compute \(D^\star(x)\) when \(p_{\mathrm{data}}(x)=0.1\) and \(p_g(x)=0.4\). Interpret the result.
\item \textbf{Interpretation.} Compare a GAN and a VAE in terms of sampling, likelihood, and inference \(z\mid x\).
\item \textbf{Misconception/debugging.} A student trains a GAN that produces one excellent-looking digit for every noise input. Name the failure mode and explain why the discriminator objective may not have enforced diversity well enough.
\end{enumerate}

\subsection{41.4 Diffusion and score-based models}

\paragraph{Guiding question.}
How can a model learn to generate data by reversing a gradual stochastic corruption process?

\paragraph{Beginner-facing intuition.}
Diffusion models turn data into noise in many small steps, then learn to reverse those steps. The forward process is fixed and easy: add Gaussian noise until the data distribution becomes nearly standard normal. The reverse process is learned: starting from noise, repeatedly denoise to produce a sample.

The score function
\[
\nabla_x\log p_t(x)
\]
points toward directions of increasing log density at noise level \(t\). Score-based models learn this vector field. In continuous time, the forward and reverse processes can be written as stochastic differential equations, so the \(dW_t\) notation from stochastic processes becomes part of generative modeling.

\paragraph{Formal definitions and notation.}
In a discrete denoising diffusion model, a common forward noising marginal is
\[
x_t=\sqrt{\overline\alpha_t}\,x_0+\sqrt{1-\overline\alpha_t}\,\epsilon,
\qquad
\epsilon\sim\mathcal N(0,I),
\]
where \(x_0\) is data and \(\overline\alpha_t\) decreases from near \(1\) toward \(0\). A neural network \(\epsilon_\theta(x_t,t)\) is trained to predict the added noise by minimizing
\[
\boxed{
\mathbb{E}_{x_0,\epsilon,t}
\left[\|\epsilon-\epsilon_\theta(x_t,t)\|_2^2\right].
}
\]
Equivalently, the network can be parameterized to estimate the score \(\nabla_{x_t}\log p_t(x_t)\), with scaling depending on the noise convention.

In continuous time, a forward diffusion has SDE form
\[
dX_t=f(X_t,t)\,dt+g(t)\,dW_t,
\]
where \(W_t\) is a Wiener process. The reverse-time SDE involves the score:
\[
dX_t=\left[f(X_t,t)-g(t)^2\nabla_x\log p_t(X_t)\right]dt+g(t)\,d\overline W_t,
\]
read backward in time. The important beginner point is that generation requires an estimate of the score field at each noise level.

\paragraph{Core mechanism: denoising a noisy Gaussian observation.}
From
\[
x_t=\sqrt{\overline\alpha_t}x_0+\sqrt{1-\overline\alpha_t}\epsilon,
\]
if a network predicts \(\widehat\epsilon=\epsilon_\theta(x_t,t)\), then an estimate of the clean data is
\[
\widehat x_0
=\frac{x_t-\sqrt{1-\overline\alpha_t}\widehat\epsilon}
{\sqrt{\overline\alpha_t}}.
\]
Thus noise prediction is directly tied to denoising. The learned reverse sampler uses many such local denoising steps, with additional variance terms depending on the chosen sampler.

\paragraph{Concrete example.}
Let the data be one-dimensional, \(x_0=4\), and \(\overline\alpha_t=0.25\). If the sampled noise is \(\epsilon=2\), then
\[
x_t=\sqrt{0.25}(4)+\sqrt{0.75}(2)
=2+1.732=3.732.
\]
If the network predicts \(\widehat\epsilon=2\), then
\[
\widehat x_0
=\frac{3.732-\sqrt{0.75}(2)}{\sqrt{0.25}}
=4.
\]
This toy calculation shows the local algebra behind denoising, though real models learn high-dimensional conditional denoisers.

\paragraph{Computational neuroscience example.}
Diffusion models use stochastic processes as generative machinery. Neural systems also exhibit noise, variability, and diffusion-like evidence accumulation in some decision models. The mathematical connection is useful, but it should not be overstated: a diffusion image generator is not a biological claim about how cortex samples images. It is a learned reverse process built from SDE or Markov-chain ideas.

\paragraph{AI/ML example.}
Modern image, audio, video, and molecular generators often use diffusion or score-based models. The denoising network may be a U-Net, transformer, or graph network depending on the data type. Classifier-free guidance and conditioning modify the reverse process so samples follow prompts, labels, or other constraints.

\paragraph{Common misconceptions and failure modes.}
The forward noising process is not learned in the simplest diffusion models; the reverse denoiser is learned. A score is a gradient of log density, not a scalar score like accuracy. SDE notation is not ordinary deterministic calculus because \(dW_t\) has Brownian scaling. More denoising steps can improve quality but increase sampling cost. A good denoising loss does not eliminate bias from the training data distribution.

\paragraph{Exercises.}
\begin{enumerate}[label=\textbf{41.4.\arabic*.}, leftmargin=*]
\item \textbf{Mechanical.} For \(x_0=1\), \(\overline\alpha_t=0.64\), and \(\epsilon=-0.5\), compute \(x_t\).
\item \textbf{Proof/derivation.} Starting from \(x_t=\sqrt{\overline\alpha_t}x_0+\sqrt{1-\overline\alpha_t}\epsilon\), solve algebraically for \(x_0\) in terms of \(x_t\) and \(\epsilon\).
\item \textbf{Numerical/computational.} If \(\overline\alpha_t=0.25\), \(x_t=0.2\), and \(\widehat\epsilon=-1\), compute \(\widehat x_0\).
\item \textbf{Interpretation.} Explain how the generative direction in diffusion differs from the inference direction in a VAE.
\item \textbf{Misconception/debugging.} A student says, ``A diffusion model learns to add noise to images.'' Correct the statement by distinguishing the fixed forward noising process from the learned reverse denoising process.
\end{enumerate}

\paragraph{Closing bridge.}
Generative deep learning turns probability into trainable tensor programs. VAEs emphasize latent-variable inference and KL bounds; flows emphasize invertible change of variables; GANs emphasize adversarial sample quality; energy models emphasize unnormalized probability and partition functions; diffusion models emphasize stochastic processes and learned scores. Chapter 42 now asks how the representations learned by such systems are shaped geometrically, how robust they are, and how cautiously we can interpret them.

\clearpage
